% Responsable: Jerónimo (relación entre los 4 modelos)
% Tesis: los 4 modelos son el MISMO problema con distinta información sobre
% quién es pobre. Esta tabla es el contrato: nadie usa otra notación.
\subsection{Notación común y jerarquía de información}\label{sec:notacion}

Los cuatro modelos comparten los datos, la medida de pobreza y la meta; lo único
que cambia entre ellos es \emph{qué sabe el planificador} sobre cada hogar en el
momento de asignar el gasto. Para que esa diferencia sea la única, todos se
escriben con la notación de la Tabla~\ref{tab:notacion}, que adapta la de
\citet[Tablas A1--A3]{hamie2024}.

\begin{table}[ht]
\centering\small
\caption{Notación común a los cuatro modelos.}\label{tab:notacion}
\begin{tabularx}{\textwidth}{@{}llX@{}}
\toprule
Símbolo & Tipo & Significado \\
\midrule
\multicolumn{3}{@{}l}{\emph{Conjuntos}}\\
$i\in\Hog$, $j\in\Ind$ & & hogares (\cHogares) e indicadores (\cIndicadores) \\
$\Hog^0\subseteq\Hog$ & & hogares pobres en la línea base (\cPobresHog) \\
$d\in\Zon$ & & zonas geográficas: cabecera, centro poblado, rural disperso \\
$t\in\Tip$ & & tipos de hogar obtenidos por \emph{clustering} (Modelo 4) \\
$\Ind=\Upriv\cup\Vpub$ & & indicadores de bien privado y de bien público (Modelo 4) \\
$g(i,j)$ & & celda de focalización del par $(i,j)$: $d_i$ en M3; $d_i$ o $t_i$ en M4 \\
\midrule
\multicolumn{3}{@{}l}{\emph{Parámetros}}\\
$w_j$, $k$ & & peso del indicador ($\sum_j w_j=1$); umbral de pobreza ($k=1/3$) \\
$M_{ij}\in\{0,1\}$ & & privación inicial del hogar $i$ en el indicador $j$ \\
$\HS_i$, $\HW_i$ & & tamaño del hogar; factor de expansión \\
$u_{ij}$ & & personas del hogar que hay que atender para quitar la privación \\
$c_{ij}$ & & costo de quitar la privación $(i,j)$ en un hogar de la muestra \\
$V_{ij}=u_{ij}\HW_i$ & & cobertura que aporta el par, en personas de la población \\
$K_{ij}=c_{ij}\HW_i$ & & gasto que exige el par, expandido a la población \\
$R_{ij}\sim U[0,1]$ & & sorteo fijo del par, independiente entre pares (M3--M4) \\
$\MPIs$, $\MPIr$ & & $\Mcero$ de la línea base (\cMceroExacto); reducción exigida \\
$B$ & & presupuesto (objetivo de equidad) \\
\midrule
\multicolumn{3}{@{}l}{\emph{Variables}}\\
$\Ndec_{ij}\in\{0,1\}$ & & privación después de intervenir ($\Ndec_{ij}\le M_{ij}$) \\
$s_i$ & & puntaje de privación final, $s_i=\sum_j w_j\Ndec_{ij}$ \\
$b1_i\in\{0,1\}$ & & $b1_i=1$ si y solo si el hogar \emph{sale} de la pobreza ($s_i<k$) \\
$\Cont_i\ge 0$ & & contribución del hogar al $\Mcero$ final \\
$\Esf_{jg}$ & & gasto en el indicador $j$ dentro de la celda $g$ \\
$\cob_{jg}\in[0,1]$ & & cobertura del par $(j,g)$: fracción de la población privada atendida \\
\bottomrule
\end{tabularx}
\end{table}

Con información perfecta, los Modelos 1 y 2 deciden $\Ndec_{ij}$ hogar por
hogar. El Modelo 3 solo decide la cobertura $\cob_{jd}$ de cada par
(indicador, zona), y es el sorteo $R_{ij}$ el que determina qué hogares la
reciben. El Modelo 4 hace lo mismo, pero los indicadores de bien privado se
focalizan por tipo de hogar $t$ en lugar de por zona. Visto así, los cuatro
modelos forman una escalera de información:
\[
\underbrace{\text{M1, M2}}_{\text{hogar}}
\;\succ\;
\underbrace{\text{M4}}_{\text{zona y tipo de hogar}}
\;\succ\;
\underbrace{\text{M3}}_{\text{zona}},
\]
y el costo adicional de bajar cada escalón para alcanzar la misma meta es lo que
llamamos el \emph{precio de la información} (Sección~\ref{sec:comparacion}).
El extremo superior de la escalera es un caso particular del inferior: si cada
hogar fuera su propia celda, el sorteo no tendría a quién más asignar la
cobertura y el Modelo 3 coincidiría con el Modelo 2. Esta propiedad, señalada por
\citet[p.~11]{hamie2024}, la verificamos numéricamente (Sección~\ref{sec:diseno}).

\nota{Jerónimo}{Martín: en el M1 tu $y_i$ tiene el mismo significado que $b1_i$
($=1$ si el hogar sale). Si en 05 usas $y_i$, dilo en una línea.}
